\chapter{Cartan 子代数、根、权与最高权}
\label{chap:roots-weights}

$SU(2)$ 只有一个彼此独立的对角生成元，因此一个磁量子数就能标记角动量多重态中的各个权态。对 $SU(3)$ 或更高秩群，生成元大多不对易，不可能同时对角化；但我们仍希望用尽可能多的一组相容量子数标记态，并知道其余生成元会把一个态送到哪里。

先选出最大的可交换生成元集合。它们的共同本征值组成“权”，其余生成元则像角动量升降算符一样移动这些本征值；移动量在伴随表示中称为“根”。根描述代数怎样作用，权描述某个具体表示中的态怎样响应。选定正根以后，单根、最高权与 Weyl 反射便把这种作用关系组织成有限的组合数据。

\section{Cartan 子代数与根空间}

上一章已经从连续群得到 Lie 代数和伴随表示。这里取有限维复半单 Lie 代数 $\mathfrak g$。伴随表示 $\operatorname{ad}_X(Y)=[X,Y]$ 把括号变成线性算符，而 Killing 形式
\begin{equation}
B(X,Y)=\operatorname{tr}(\operatorname{ad}_X\operatorname{ad}_Y)
\label{eq:killing-form}
\end{equation}
在半单 Lie 代数上非退化。对紧半单实形式，$-B$ 正定，因而可以比较根空间中的长度和夹角。

\begin{definition}[Cartan 子代数]
Cartan 子代数 $\mathfrak h$ 是 $\mathfrak g$ 中最大的可交换半单子代数，其维数 $r$ 称为 $\mathfrak g$ 的秩。复化后，伴随表示可在 $\mathfrak h$ 上同时对角化，按根空间分解为
\[
\mathfrak g_{\mathbb C}=\mathfrak h_{\mathbb C}
\oplus\bigoplus_{\alpha\in\Phi}\mathfrak g_\alpha,
\qquad [H,E_\alpha]=\alpha(H)E_\alpha .
\]
$\alpha\in\mathfrak h^*$ 称为根，$\Phi$ 为根系。
\end{definition}

设 $V$ 是 $\mathfrak g$ 的一个有限维表示，表示映射记为 $\rho$。由于 $\mathfrak h$ 中的生成元彼此对易，可以同时对角化。若非零向量 $v_\mu\in V$ 满足
\begin{equation}
\rho(H)v_\mu=\mu(H)v_\mu,
\qquad H\in\mathfrak h,
\label{eq:weight-definition}
\end{equation}
则线性泛函 $\mu\in\mathfrak h^*$ 称为一个权。现在让根向量 $E_\alpha$ 作用在 $v_\mu$ 上：
\begin{align*}
\rho(H)\rho(E_\alpha)v_\mu
&=\rho(E_\alpha)\rho(H)v_\mu
 +\rho([H,E_\alpha])v_\mu\\
&=[\mu(H)+\alpha(H)]\rho(E_\alpha)v_\mu.
\end{align*}
只要结果不为零，$E_\alpha v_\mu$ 就是权为 $\mu+\alpha$ 的新态。这个计算给出了根图与权图的直接联系：根向量沿根的方向搬运权态。

在 $\mathfrak{su}(2)$ 中，可取 Cartan 生成元 $J_z$，而 $J_\pm=J_x\pm\ii J_y$ 满足 $[J_z,J_\pm]=\pm\hbar J_\pm$。因此权就是 $m\hbar$，两个根对应步长 $\pm\hbar$。更高秩代数只是同时拥有多个对角生成元，权和根因而成为多分量向量。

用 Killing 型把 $\mathfrak h$ 与 $\mathfrak h^*$ 识别后，根的内积记为
\begin{equation}
(\alpha,\beta)=B(H_\alpha,H_\beta).
\label{eq:root-inner-product}
\end{equation}

\begin{derivation}[在 $su(3)$ 上验证 Killing 型非退化]
目标：直接验算 $B(X,Y)=\tr(\operatorname{ad}_X\operatorname{ad}_Y)$ 在 $\mathfrak{su}(3)$ 上非退化。取 $\mathfrak h$ 的基 $H_1,H_2$，使其满足 $\alpha_i(H_j)=\delta_{ij}$；$su(3)$ 的六个根为 $\pm\alpha_1,\pm\alpha_2,\pm(\alpha_1+\alpha_2)$。对 $H\in\mathfrak h$，$\operatorname{ad}_H$ 在根空间 $\mathfrak g_\alpha$ 上作用为标量 $\alpha(H)$，而在 $\mathfrak h$ 上作用为零，故
\[
B(H,H')=\sum_{\alpha\in\Phi}\alpha(H)\alpha(H')
=2\sum_{\alpha>0}\alpha(H)\alpha(H').
\]
记 $H=(h_1,h_2)$（即 $\alpha_1(H)=h_1,\alpha_2(H)=h_2$）、$H'=(h_1',h_2')$，三个正根贡献
\[
B(H,H')=2\bigl[h_1h_1'+h_2h_2'+(h_1+h_2)(h_1'+h_2')\bigr]
=\begin{pmatrix}h_1&h_2\end{pmatrix}
\begin{pmatrix}4&2\\2&4\end{pmatrix}
\begin{pmatrix}h_1'\\h_2'\end{pmatrix}.
\]
该矩阵行列式 $4\cdot4-2\cdot2=12>0$，本征值为 $6,2$，故 $B$ 在 $\mathfrak h$ 上正定、非退化。半单性又保证 $B$ 在 $\mathfrak h$ 与根空间 $\mathfrak g_\alpha$ 之间的分块为零（迹只在同一权空间内贡献），故 $B$ 在整个 $su(3)$ 上非退化。这正是可用式 \eqnrefs{eq:root-inner-product} 与式 \eqnrefs{eq:cartan-matrix} 在 $\mathfrak h^*$ 上定义根内积与 Cartan 矩阵的前提。
\end{derivation}

\section{正根、单根与最高权}

根总是成对出现为 $\pm\alpha$，因此必须先选定“向上”的方向。选择正根后，每个正根都可由简单根 $\alpha_i$ 的非负整数线性组合表示。Cartan 矩阵
\begin{equation}
A_{ij}=\frac{2(\alpha_i,\alpha_j)}{(\alpha_j,\alpha_j)}
\label{eq:cartan-matrix}
\end{equation}
把简单根的相对长度和夹角编码为整数矩阵。有限维复半单 Lie 代数的 Dynkin 图正是由这个矩阵分类的。

\begin{definition}[单根、基本权与 Dynkin 标号]
选定正根后，每个正根都可写成单根 $\alpha_i$（$i=1,\dots,r$）的非负整线性组合。基本权 $\omega_i\in\mathfrak h^*$ 由
\[
\frac{2(\omega_i,\alpha_j)}{(\alpha_j,\alpha_j)}=\delta_{ij}
\]
定义。任意有限维表示中 $\mathfrak h$ 的共同本征值 $\mu\in\mathfrak h^*$ 称为权；支配权 $\lambda=\sum_{i=1}^r a_i\omega_i$（$a_i\in\mathbb Z_{\ge0}$）唯一标记一个不可约有限维表示，非负整数 $a_i$ 即 Dynkin 标号。这个结论把不可约表示的分类化为一个离散格点问题，而降算符 $E_{-\alpha_i}$ 从最高权逐层生成全部权。
\end{definition}

若权 $\mu$ 的重数为 $m(\mu)$，Freudenthal 递推式（式 \eqnrefs{eq:freudenthal}）给出
\begin{equation}
\bigl[(\lambda+\rho,\lambda+\rho)-(\mu+\rho,\mu+\rho)\bigr]m(\mu)
=2\sum_{\alpha>0}\sum_{k\ge1}(\mu+k\alpha,\alpha)m(\mu+k\alpha),
\label{eq:freudenthal}
\end{equation}
其中 $\rho=\frac12\sum_{\alpha>0}\alpha$，右侧只含比 $\mu$ 更高的权，因此递推从 $m(\lambda)=1$ 开始闭合。

紧群的有限维幺正不可约表示还满足 Peter--Weyl 定理：矩阵元在 $L^2(G)$ 中构成正交完备系，
\[
L^2(G)\cong\widehat{\bigoplus}_{\pi\in\widehat G}V_\pi\otimes V_\pi^* .
\]
这既是有限群正则表示分解的连续版本，也是紧群 Fourier 分析的基础。非紧群一般需要直积分与 Plancherel 测度，不能直接照搬离散直和。

\begin{derivation}[Freudenthal 递推的来源]
二次 Casimir $C_2=\sum_aT^aT^a$ 在最高权表示 $V_\lambda$ 上取标量 $(\lambda,\lambda+2\rho)$。把生成元分成 Cartan 基与根向量，并在权空间 $V_\mu$ 上取迹，根向量只连接 $\mu$ 与 $\mu\pm\alpha$。用 $[E_\alpha,E_{-\alpha}]=H_\alpha$ 把成对项移到更高权空间，便得到
\[
\bigl[(\lambda+\rho)^2-(\mu+\rho)^2\bigr]m(\mu)
=2\sum_{\alpha>0}\sum_{k\ge1}(\mu+k\alpha,\alpha)m(\mu+k\alpha).
\]
递推之所以终止，是因为有限维最高权表示中，沿每条正根方向只有有限长的权串。
\end{derivation}

\begin{example}[Freudenthal 递推在 $SU(3)$ 八重态 $(1,1)$ 上走一步]
取最高权表示 $(1,1)$，其最高权 $\lambda=\omega_1+\omega_2=\rho$，已知 $m(\lambda)=1$。递推第一个要算的是次高权 $\mu=\lambda-\alpha_1=\rho-\alpha_1=\alpha_2$。用上面的归一化 $(\alpha_i,\alpha_i)=1$、$(\rho,\alpha_1)=\tfrac12$。

\emph{左边}：$\lambda+\rho=2\rho$，$\mu+\rho=2\rho-\alpha_1$，故
\[
(\lambda+\rho,\lambda+\rho)-(\mu+\rho,\mu+\rho)
=4(\rho,\alpha_1)-(\alpha_1,\alpha_1)
=4\cdot\tfrac12-1=1.
\]

\emph{右边}：$\mu=\alpha_2$。逐根检查 $\mu+k\alpha$ 是否落在已知权上：对 $\alpha=\alpha_1$、$k=1$ 得 $\mu+\alpha_1=\rho=\lambda$，$m=1$，贡献 $(\mu+\alpha_1,\alpha_1)m(\lambda)=(\rho,\alpha_1)=\tfrac12$；对 $\alpha=\alpha_2$ 与 $\alpha_1+\alpha_2$，$\mu+k\alpha$ 落到 $\alpha_1+\alpha_2$ 之外，不在八重态权图上，$m=0$。故右边
\[
2\sum_{\alpha>0}\sum_{k\ge1}(\mu+k\alpha,\alpha)m(\mu+k\alpha)
=2\cdot\tfrac12=1.
\]

两边相等，$1\cdot m(\mu)=1$，得 $m(\alpha_2)=1$，与八重态（伴随表示）每个根空间一重的事实一致。其余权可同法递推闭合。
\end{example}


\section{Weyl 群与特征标公式}

对每个根 $\alpha$，根空间中的反射为
\[
s_\alpha(\lambda)=\lambda-
\frac{2(\lambda,\alpha)}{(\alpha,\alpha)}\alpha.
\]
由简单根反射生成的有限群称为 Weyl 群。它保持根系和内积，并把权空间分成 Weyl 室；每条 Weyl 轨道恰有一个支配权落在闭基本室中。

若 $\lambda$ 是支配整权，Weyl 特征标公式（式 \eqnrefs{eq:weyl-character}）为
\begin{equation}
\chi_\lambda(\ee^H)=
\frac{\sum_{w\in W}\det(w)\,
\ee^{(w(\lambda+\rho),H)}}
{\sum_{w\in W}\det(w)\,\ee^{(w\rho,H)}}.
\label{eq:weyl-character}
\end{equation}
令 $H\to0$ 时分子分母同时为零，须用行列式展开取最低非零阶。把分子按 $W$ 中元素排成行列式：$\sum_{w\in W}\det(w)\,\ee^{(w(\lambda+\rho),H)}$ 是关于变量 $x_\alpha=\ee^{(\alpha,H)}$ 的斜反对称多项式，在 $H\to0$（即 $x_\alpha\to1$）时其最低非零阶就是 Vandermonde 型
\[
\prod_{\alpha>0}\bigl(\ee^{(\lambda+\rho,H)_\alpha/2}-\ee^{-(\lambda+\rho,H)_\alpha/2}\bigr)
\]
中首次展开的线性项，比例系数恰为 $\prod_{\alpha>0}(\lambda+\rho,\alpha)$；分母同法给出 $\prod_{\alpha>0}(\rho,\alpha)$。二者相除、公共的二阶小因子抵消，得到 Weyl 维数公式（式 \eqnrefs{eq:weyl-dimension}）
\begin{equation}
\dim V_\lambda=
\prod_{\alpha>0}
\frac{(\lambda+\rho,\alpha)}{(\rho,\alpha)}.
\label{eq:weyl-dimension}
\end{equation}
以 $(p,q)=(1,1)$ 检验：$\lambda+\rho=2\rho$，三个正根的比值分别为 $2,2,2$，乘积 $8$，与八重态维数一致。
它给出维数，特征标的乘积则给出张量积分解。例如
\[
\mathbf8\otimes\mathbf8
=\mathbf1\oplus\mathbf8_s\oplus\mathbf8_a
\oplus\mathbf{10}\oplus\overline{\mathbf{10}}\oplus\mathbf{27},
\]
而\eqnrefs{eq:baryon-decomposition}的右边 $\mathbf{10}\oplus\mathbf8\oplus\mathbf8\oplus\mathbf1$ 正是重子十重态、两个八重态与单态。
